\section{Related Work}\label{sec:related}
\textbf{Visual robustness of VLA policies.} LIBERO-Plus~\cite{fei2026liberoplus} and related evaluations~\cite{guo2026robustvla,wang2025adversarial} report substantial success losses under visual, multi-modal, and adversarial perturbations. Robustness methods edit task-irrelevant image regions during deployment~\cite{byovla} or train for consistent actions under task-preserving input changes~\cite{luo2026rovla,guo2026robustvla,peng2026pair}. Task success, their common measure, shows that robustness is lost but not what a perturbation changes inside the policy, or whether larger action changes mean larger losses. We measure both.

\textbf{Mechanistic analysis and internal interventions.} Prior analyses trace representations and knock out attention~\cite{vlatrace}, patch activations~\cite{meng2022locating} to locate language grounding~\cite{wulff2026language}, and relate interventional visual attribution to generalization~\cite{zhang2026embodied}. Others find features and directions that steer robot behavior~\cite{haon2025,sae,observing,grant2026}, or steer activations toward subspaces fitted on successful and failed rollouts, which transfer across tasks with similar failure modes~\cite{miao2026coast}. We instead study the feature changes that visual perturbations cause, which any internal correction of visual robustness must act on. We fit directions to these changes on clean--perturbed pairs matched in state and policy noise, without outcome labels, hold out tasks and families together, and test the directions by a subspace form of activation patching during the policy's own generation.

\textbf{Evaluating robot behavior beyond task success.} Similar success rates can accompany different smoothness, efficiency, and gripper behavior in successful trajectories~\cite{higham2026success}. Matched executions under individual and combined perturbations reveal opposing outcome transitions that cancel in aggregate success~\cite{sawada2026paired}, and out-of-distribution inputs alone do not predict failure~\cite{zhu2026vlascope}. Run-to-run variation is a known hazard when comparing learned policies~\cite{agarwal2021precipice}, and samples that a model update turns from correct to wrong are tracked as negative flips in classification~\cite{yan2021positive}. We bring both to perturbed VLA executions: the variation between two clean runs that differ only in policy noise is the scale, and lost and gained successes are tracked under single perturbations.
